\chapter*{Abstract}
\addcontentsline{toc}{chapter}{Abstract}
\markright{Abstract}

Fine-tuning one pre-trained language model per task gives strong task-specific performance, but every new task adds a complete model to train, store and serve, and nothing learned for one task reaches the others. This study examines whether several fine-tuned models can become a single multi-task model whose weights adapt to the input. We adapt Weight-Ensembling Mixture of Experts (WEMoE), previously evaluated on vision Transformers, to BERT and five binary tasks of the GLUE benchmark: MRPC, RTE, CoLA, SST-2 and QNLI. Weights outside the MLP blocks are merged statically with Task Arithmetic, while the MLP weights of every encoder layer are rebuilt for each input as the pre-trained weights plus a router-weighted sum of frozen task vectors. Only the routers are trained, by Few-Shot Fine-Tuning (FSFT) on labelled examples that the experts never saw. All model selection uses a development split, and scores are reported on the official validation splits. With 256 examples per task, WEMoE reaches an average score of $69.9\pm0.1$, against 57.4 for Task Arithmetic, 78.2 for the experts and 45.8 for a majority-class predictor, which is 59.1\% of the gain between the majority class and the experts. Coefficients that do not depend on the input reach only $58.4\pm0.0$. Upscaling only the first MLP layer gives a similar average ($69.1\pm0.1$) with 47.2\% of the storage of five experts rather than 73.1\%, but the two variants are weak on different tasks. WEMoE matches the experts on no task and is several times slower in batches. Limitations include a single set of experts, small validation sets for two tasks, and development data that exceeds the few-shot budget.

\vspace{3mm}
\noindent\textbf{Keywords:} model consolidation; weight ensembling; Mixture of Experts; BERT; multi-task learning; few-shot fine-tuning.
